% Part: first-order-logic
% Chapter: syntax-and-semantics
% Section: formation-sequences

\documentclass[../../../include/open-logic-section]{subfiles}

\begin{document}

\olfileid[pt]{pl}{syn}{fseq}

\olsection{Sequências de Formação}

Definir !!{formula}s por meio de uma definição indutiva, juntamente com a
técnica complementar de demonstrar propriedades de !!{formula}s por indução,
é uma abordagem elegante e eficiente. Contudo, também pode ser útil considerar
uma abordagem mais ascendente e passo a passo para a construção de
!!{formula}s, como fazemos aqui usando a noção de \emph{sequência de formação}.

\begin{defn}[Sequências de formação para fórmulas]
\ollabel{defn:fseq-frm}
Uma sequência finita $\tuple{!A_0,\dotsc,!A_n}$ de cadeias de símbolos da
linguagem~$\Lang L_0$ é uma \emph{sequência de formação} para $!A$ se
$!A \ident !A_n$ e, para todo $i \leq n$, ou $!A_i$ é uma fórmula atômica,
ou existem $j,k < i$ tais que vale uma das seguintes condições:
\begin{enumerate}
    \tagitem{prvNot}{$!A_i \ident \lnot !A_j$.}{}%
    \tagitem{prvAnd}{$!A_i \ident (!A_j \land !A_k)$.}{}%
    \tagitem{prvOr}{$!A_i \ident (!A_j \lor !A_k)$.}{}%
    \tagitem{prvIf}{$!A_i \ident (!A_j \lif !A_k)$.}{}%
    \tagitem{prvIff}{$!A_i \ident (!A_j \liff !A_k)$.}{}%
\end{enumerate}
\end{defn}

\begin{ex}
\[
\tuple{
    \Obj p_0,
    \Obj p_1,
    (\Obj p_1 \land \Obj p_0),
    \lnot (\Obj p_1 \land \Obj p_0)
}
\]
é uma sequência de formação de
$\lnot (\Obj p_1 \land \Obj p_0)$, assim como
\[
\tuple{
    \Obj p_0,
    \Obj p_1,
    \Obj p_0,
    (\Obj p_1 \land \Obj p_0),
    (\Obj p_0 \lif \Obj p_1),
    \lnot (\Obj p_1 \land \Obj p_0)
}.
\]
%
Como se pode ver no segundo exemplo, sequências de formação podem conter
`lixo': fórmulas redundantes ou que não contribuem para a construção.
\end{ex}

\begin{prop}
\ollabel{prop:formed}
Toda !!{formula}~$!A$ em~$\Frm[L_0]$ tem uma sequência de formação.
\end{prop}

\begin{proof}
Suponha que $!A$ seja atômica. Então a sequência~$\tuple{!A}$ é uma
sequência de formação para~$!A$.
%
Agora suponha que $!B$ e~$!C$ tenham, respectivamente, as sequências de
formação $\tuple{!B_0,\dotsc,!B_n}$ e $\tuple{!C_0,\dotsc,!C_m}$.
%
\begin{enumerate}
  \tagitem{prvNot}{Se $!A \ident \lnot !B$,
      então $\tuple{!B_0,\dotsc,!B_n,\lnot !B_n}$
      é uma sequência de formação para~$!A$.}{}
  \tagitem{prvAnd}{Se $!A \ident (!B \land !C)$,
      então $\tuple{!B_0,\dotsc,!B_n,!C_0,\dotsc,!C_m,(!B_n \land !C_m)}$
      é uma sequência de formação para~$!A$.}{}
  \tagitem{prvOr}{Se $!A \ident (!B \lor !C)$,
      então $\tuple{!B_0,\dotsc,!B_n,!C_0,\dotsc,!C_m,(!B_n \lor !C_m)}$
      é uma sequência de formação para~$!A$.}{}
  \tagitem{prvIf}{Se $!A \ident (!B \lif !C)$,
      então $\tuple{!B_0,\dotsc,!B_n,!C_0,\dotsc,!C_m,(!B_n \lif !C_m)}$
      é uma sequência de formação para~$!A$.}{}
  \tagitem{prvIff}{Se $!A \ident (!B \liff !C)$,
      então $\tuple{!B_0,\dotsc,!B_n,!C_0,\dotsc,!C_m,(!B_n \liff !C_m)}$
      é uma sequência de formação para~$!A$.}{}
\end{enumerate}
Pelo princípio de indução sobre !!{formula}s, toda !!{formula} tem uma
sequência de formação.
\end{proof}

Também podemos demonstrar a recíproca. Isso é importante porque mostra que
nossas duas maneiras de definir fórmulas são equivalentes: elas produzem os
mesmos resultados. Também significa que podemos demonstrar teoremas sobre
fórmulas usando indução ordinária sobre o comprimento das sequências de
formação.

\begin{lem}
\ollabel{lem:fseq-init}
Suponha que $\tuple{!A_0,\dotsc,!A_n}$ seja uma sequência de formação
para~$!A_n$ e que $k \leq n$. Então $\tuple{!A_0,\dotsc,!A_k}$ é uma
sequência de formação para~$!A_k$.
\end{lem}

\begin{proof}
Exercício.
\end{proof}

\begin{thm}
\ollabel{thm:fseq-frm-equiv}
$\Frm[L_0]$ é o conjunto de todas as cadeias de símbolos da
linguagem~$\Lang L_0$ que têm uma sequência de formação.
\end{thm}

\begin{proof}
Seja $F$ o conjunto de todas as cadeias de símbolos da
linguagem~$\Lang L_0$ que têm uma sequência de formação.
Vimos em \olref[pl][syn][fseq]{prop:formed} que
$\Frm[L_0] \subseteq F$; agora demonstramos a recíproca.

Suponha que $!A$ tenha uma sequência de formação
$\tuple{!A_0,\dotsc,!A_n}$. Demonstramos que $!A \in \Frm[L_0]$ por
indução forte sobre~$n$. Nossa hipótese de indução é que toda cadeia de
símbolos com uma sequência de formação de comprimento $m < n$ pertence
a~$\Frm[L_0]$. Pela definição de sequência de formação, ou $!A_n$ é
atômica, ou devem existir $j,k < n$ tais que vale um dos seguintes casos:
\begin{enumerate}
    \tagitem{prvNot}{$!A_n \ident \lnot !A_j$.}{}%
    \tagitem{prvAnd}{$!A_n \ident (!A_j \land !A_k)$.}{}%
    \tagitem{prvOr}{$!A_n \ident (!A_j \lor !A_k)$.}{}%
    \tagitem{prvIf}{$!A_n \ident (!A_j \lif !A_k)$.}{}%
    \tagitem{prvIff}{$!A_n \ident (!A_j \liff !A_k)$.}{}%
\end{enumerate}
Agora raciocinamos por casos. Se $!A_n$ é atômica, então
$!A_n \in \Frm[L_0]$. Suponha, em vez disso, que $!A \equiv
(!A_j \land !A_k)$. Por \olref[pl][syn][fseq]{lem:fseq-init},
$\tuple{!A_0,\dotsc,!A_j}$ e $\tuple{!A_0,\dotsc,!A_k}$ são,
respectivamente, sequências de formação para $!A_j$ e~$!A_k$. Como são
subsequências iniciais próprias da sequência de formação para~$!A$, ambas
têm comprimento menor que~$n$. Portanto, pela hipótese de indução, $!A_j$
e~$!A_k$ pertencem a~$\Frm[L_0]$; assim, pela definição de !!a{formula},
$(!A_j \land !A_k)$ também pertence. Os outros casos decorrem de raciocínios
paralelos.
\end{proof}

\end{document}
